\documentclass[11pt,a4paper]{article}
\usepackage[utf8]{inputenc}
\usepackage[vietnamese]{babel}
\usepackage{amsmath,amssymb}
\usepackage{geometry}
\geometry{a4paper, total={180mm,260mm}, left=15mm, right=15mm, top=18mm, bottom=18mm}
\usepackage{tikz}
\usetikzlibrary{shapes.geometric, arrows.meta, positioning, calc}
\usepackage{tcolorbox}
\tcbuselibrary{skins,breakable}
\usepackage{xcolor}
\usepackage{multicol}
\usepackage{array}
\usepackage{fontawesome}
\usepackage{fancyhdr}

\pagestyle{fancy}
\fancyhf{}
\renewcommand{\headrulewidth}{0.5pt}
\renewcommand{\footrulewidth}{0.5pt}
\fancyhead[L]{\small\bfseries\color{teal}\faBook\ SỔ TAY KIẾN THỨC TOÁN 3}
\fancyhead[R]{\small\bfseries\color{teal}Lớp Toán Thầy Lĩnh -- loptoanthaylinh.vn}
\fancyfoot[C]{\small\color{gray}-- Sổ tay kiến thức: Bài 74. Khả năng xảy ra của một sự kiện --}

\definecolor{primary}{HTML}{005F73}
\definecolor{secondary}{HTML}{0A9396}
\definecolor{accent}{HTML}{EE9B00}
\definecolor{danger}{HTML}{AE2012}
\definecolor{boxbg}{HTML}{F4F6F8}

\begin{document}

% --- BANNER TIÊU ĐỀ ---
\begin{tcolorbox}[
  enhanced,
  colback=white,
  colframe=primary,
  arc=3mm,
  boxrule=1.2pt,
  left=4mm, right=4mm, top=2.5mm, bottom=2.5mm,
  drop shadow=black!10!white
]
  \centering
  {\Large\bfseries\color{primary}\faBookmark\ BÀI 74: KHẢ NĂNG XẢY RA CỦA MỘT SỰ KIỆN}\\[1.5mm]
  {\small\itshape\color{black!75}Ba mức độ xác suất sơ khai: Chắc chắn, Có thể, Không thể qua các tình huống lấy bóng và tung xúc xắc trên 1 trang A4}
\end{tcolorbox}

\vspace{-2mm}


\noindent
\begin{minipage}[t]{0.49\textwidth}\vspace{0pt}
\begin{tcolorbox}[
  enhanced,
  colback=blue!4!white,
  colframe=blue!75!black,
  arc=2mm, boxrule=0.8pt,
  left=2mm, right=2mm, top=1.8mm, bottom=1.8mm,
  title={\footnotesize\bfseries\color{white}\faKey\ 1. Ba mức độ khả năng của sự kiện}
]
{\footnotesize
\begin{itemize}\setlength{\itemsep}{1.5pt}\setlength{\leftskip}{-2.5mm}
  \item Khi thực hiện một hành động (như rút thẻ, lấy bóng, tung đồng xu), một sự kiện có thể có $3$ mức độ khả năng:
  \item \textbf{1. Chắc chắn:} Sự kiện \textbf{luôn luôn xảy ra} trong mọi trường hợp (không thể nào khác).
  \item \textbf{2. Có thể:} Sự kiện \textbf{có lúc xảy ra, có lúc không} xảy ra (phụ thuộc vào may rủi).
  \item \textbf{3. Không thể:} Sự kiện \textbf{hoàn toàn không bao giờ xảy ra} (trái với điều kiện thực tế).
\end{itemize}
}
\end{tcolorbox}
\end{minipage}
\hfill
\begin{minipage}[t]{0.49\textwidth}\vspace{0pt}
\begin{tcolorbox}[
  enhanced,
  colback=teal!4!white,
  colframe=teal!75!black,
  arc=2mm, boxrule=0.8pt,
  left=2mm, right=2mm, top=1.8mm, bottom=1.8mm,
  title={\footnotesize\bfseries\color{white}\faSliders\ 2. Mô hình phân tích thực nghiệm}
]
{\footnotesize
\begin{itemize}\setlength{\itemsep}{1.5pt}\setlength{\leftskip}{-2.5mm}
  \item \textbf{Hộp chỉ toàn bóng đỏ:}\\
  -- Lấy được bóng đỏ: \textbf{Chắc chắn}.\\
  -- Lấy được bóng xanh: \textbf{Không thể}.
  \item \textbf{Hộp có cả bóng đỏ và bóng xanh:}\\
  -- Lấy được bóng đỏ: \textbf{Có thể}.\\
  -- Lấy được bóng xanh: \textbf{Có thể}.\\
  -- Lấy được bóng vàng: \textbf{Không thể}.
  \item \textbf{Gieo xúc xắc 6 mặt (1 đến 6 chấm):}\\
  Xuất hiện mặt $7$ chấm: \textbf{Không thể}.
\end{itemize}
}
\end{tcolorbox}
\end{minipage}

\vspace{-2.5mm}
% --- KHỐI 3: SƠ ĐỒ TƯ DUY TIKZ ---
\noindent
\begin{tcolorbox}[
  enhanced,
  colback=orange!4!white,
  colframe=orange!80!black,
  arc=2mm, boxrule=0.8pt,
  left=2mm, right=2mm, top=1.5mm, bottom=1.5mm,
  title={\footnotesize\bfseries\color{white}\faSitemap\ SƠ ĐỒ TƯ DUY: BA CẤP ĐỘ KHẢ NĂNG CỦA SỰ KIỆN}
]
\centering
\begin{tikzpicture}[
  box/.style={draw, rounded corners=2pt, align=center, font=\scriptsize\bfseries, inner sep=2.5pt, minimum height=0.75cm},
  arr/.style={-{Stealth[length=2mm]}, thick}
]
  \node[box, fill=green!20, draw=green!70!black, text width=3.2cm] (k1) at (0, 0) {CHẮC CHẮN\\$100\%$ xảy ra\\$100\%$ đối tượng thỏa mãn};
  \node[box, fill=blue!15, draw=blue!80!black, text width=3.2cm] (k2) at (4.6, 0) {CÓ THỂ\\Có thể xảy ra / không\\Có một số đối tượng thỏa};
  \node[box, fill=red!15, draw=red!70!black, text width=3.2cm] (k3) at (9.2, 0) {KHÔNG THỂ\\$0\%$ xảy ra\\Không có đối tượng nào};

  \draw[arr, teal!80!black] (k1) -- (k2);
  \draw[arr, orange!80!black] (k2) -- (k3);
\end{tikzpicture}
\end{tcolorbox}

\vspace{-2.5mm}
% --- KHỐI 4: VÍ DỤ & MẸO TRÁNH BẪY ---
\noindent
\begin{minipage}[t]{0.49\textwidth}\vspace{0pt}
\begin{tcolorbox}[
  enhanced,
  colback=green!4!white,
  colframe=green!60!black,
  arc=2mm, boxrule=0.8pt,
  left=2mm, right=2mm, top=1.8mm, bottom=1.8mm,
  title={\footnotesize\bfseries\color{white}\faCheckCircle\ Ví dụ thực hành chuẩn mực}
]
{\footnotesize
\textbf{Trong túi có 3 viên bi đỏ và 2 viên bi xanh. Không nhìn vào túi, lấy ra 1 viên bi:}\\
-- Sự kiện "Lấy được viên bi đỏ": \textbf{Có thể}.\\
-- Sự kiện "Lấy được viên bi xanh": \textbf{Có thể}.\\
-- Sự kiện "Lấy được viên bi vàng": \textbf{Không thể}.\\
-- Sự kiện "Lấy được viên bi đỏ hoặc xanh": \textbf{Chắc chắn}.
}
\end{tcolorbox}
\end{minipage}
\hfill
\begin{minipage}[t]{0.49\textwidth}\vspace{0pt}
\begin{tcolorbox}[
  enhanced,
  colback=red!4!white,
  colframe=red!70!black,
  arc=2mm, boxrule=0.8pt,
  left=2mm, right=2mm, top=1.8mm, bottom=1.8mm,
  title={\footnotesize\bfseries\color{white}\faExclamationTriangle\ Mẹo tránh bẫy trắc nghiệm}
]
{\footnotesize
\begin{itemize}\setlength{\itemsep}{1.5pt}\setlength{\leftskip}{-2.5mm}
  \item \textbf{Bẫy nhầm giữa "Có thể" và "Chắc chắn":}\\
  Nếu trong túi có $99$ viên đỏ và $1$ viên xanh, việc lấy được viên đỏ vẫn chỉ là \textbf{"Có thể"} (chưa phải chắc chắn $100\%$, vì vẫn có xác suất rút phải viên xanh)!
  \item Chỉ dùng từ "Chắc chắn" khi trong túi \textbf{chỉ có duy nhất} màu đỏ.
\end{itemize}
}
\end{tcolorbox}
\end{minipage}


\end{document}
